\documentclass[11pt]{article}
\usepackage{amsmath,amssymb,amsthm,booktabs,mathtools}
\usepackage[margin=1in]{geometry}
\usepackage{microtype}
\usepackage{url}
\newtheorem{theorem}{Theorem}[section]
\newtheorem{lemma}[theorem]{Lemma}
\newtheorem{corollary}[theorem]{Corollary}
\newtheorem{remark}[theorem]{Remark}
\newcommand{\dd}{\,\mathrm d}
\newcommand{\one}{\mathbf 1}
\newcommand{\ip}[2]{\left\langle #1,#2\right\rangle}
\title{Atlas 178: A Half-Degree Weighted \(K_4\) Ratio}
\author{}
\date{}
\begin{document}
\maketitle

\begin{abstract}
Graph Atlas 178 is formed from two copies of \(K_4\) sharing a common
triangle, with one pendant leaf at one of the two page vertices.  We prove
for every graphon \(W\) of edge density \(p\in[2/3,1]\) that
\[
 t(H_{178},W)\geq p^2(2p-1)(3p-2)^2,
\]
its exact chromatic target.  Cauchy--Schwarz on the two pages reduces the
problem to a \(K_4\) whose page vertex carries the fractional degree weight
\(d^{1/2}\).  Rooted-link Goodman and two sharp supporting planes in the
degree, its one-step average, and the rooted triangle prove the required
weighted ratio.  The supporting planes are certified on the complete
rooted feasible region by exact rational conditional Bernstein
certificates.  The proof is for arbitrary graphons and uses no finite-step
reduction.
\end{abstract}

\section{The graph and its target}

Let \(H_{178}\) have a triangle on vertices \(s_1,s_2,s_3\), two page
vertices \(u,v\), each adjacent to all three \(s_i\), and a leaf adjacent to
\(u\).  This graph is isomorphic to NetworkX Graph Atlas 178 and has graph6
code \verb|En}?|.  A proper colouring first chooses ordered distinct
colours on the spine, then a colour outside those three for each page, and
finally a colour different from the colour of \(u\).  Hence
\[
 P_{178}(z)=z(z-1)^2(z-2)(z-3)^2.
\]
Writing
\[
 c_3=2p-1,\qquad c_4=3p-2,
\]
gives
\begin{equation}
 (1-p)^6P_{178}\!\left(\frac1{1-p}\right)
 =p^2c_3c_4^2.
 \label{eq:target}
\end{equation}
The required density interval is \(p\geq2/3\).

\section{Rooted data}

Write \(T_W\) for the integral operator of \(W\), and put
\begin{align}
 p&=\int_{\Omega^2}W(x,y)\dd\mu(x)\dd\mu(y),
 &d(x)&=(T_W\one)(x),\notag\\
 a(x)&=(T_Wd)(x),
 &\tau(x)&=\int_{\Omega^2}
 W(x,y)W(x,z)W(y,z)\dd\mu(y)\dd\mu(z).
 \label{eq:rooted}
\end{align}
Let
\[
 T=\int_\Omega\tau(x)\dd\mu(x)=t(K_3,W),
\qquad
 I=\int_\Omega d(x)^{1/2}\tau(x)\dd\mu(x).
\]
The standard rooted Goodman argument gives
\begin{equation}
 T\geq p(2p-1)=pc_3.
 \label{eq:goodman}
\end{equation}

We use the following complete scalar region.  It is repeated here so the
supporting-plane certificates are self-contained.

\begin{lemma}[Rooted feasible region]\label{lem:feasible}
For almost every \(x\), the tuple
\((d,a,t)=(d(x),a(x),\tau(x))\) satisfies
\begin{equation}
\begin{gathered}
 0\leq d\leq1,\qquad
 \max\{0,d+p-1\}\leq a\leq d,\\
 \max\{0,2a-p,a-d(1-d)\}\leq t\leq\min\{a,d^2\}.
\end{gathered}
\label{eq:feasible}
\end{equation}
Moreover,
\begin{equation}
 \int d=p,\qquad
 \int a=\int d^2,\qquad
 \int t=T.
 \label{eq:means}
\end{equation}
\end{lemma}

\begin{proof}
Deleting factors gives \(a\leq d\), \(t\leq a\), and \(t\leq d^2\).
Writing \(U=1-W\) gives
\[
 a(x)=p-\int U(x,y)d(y)\dd\mu(y)\geq d(x)+p-1.
\]
The two inequalities
\[
 W(x,y)W(x,z)\geq W(x,y)+W(x,z)-1
\]
and
\[
 W(x,z)W(y,z)\geq W(x,z)+W(y,z)-1
\]
give, after multiplication by the remaining edge factor and integration,
\(t\geq2a-p\) and \(t\geq a-d(1-d)\), respectively.  Nonnegativity supplies
the remaining bounds.  The identities in \eqref{eq:means} are symmetry and
Fubini.
\end{proof}

\section{Two half-degree supporting planes}

The next lemma is the only scalar input particular to this proof.

\begin{lemma}[Half-degree rooted planes]\label{lem:planes}
Let \(p\in[2/3,1]\), set \(r=\sqrt p\), and let
\((d,a,t)\) satisfy \eqref{eq:feasible}.  Then
\begin{align}
 (2\sqrt d-r)t
 &\geq
 prc_3+\frac32r(4p-1)(d-p)
       +\frac32r(a-d^2),
 \label{eq:first-plane}\\
 \frac{t}{\sqrt d}\bigl(2t-d^2-c_4d\bigr)
 &\geq
 3prc_3(d-p)+2rc_3(a-d^2).
 \label{eq:second-plane}
\end{align}
In \eqref{eq:second-plane}, the first term is defined to be zero at
\(d=0\); feasibility then forces \(a=t=0\).
\end{lemma}

The exact proof of Lemma~\ref{lem:planes} is given in
Section~\ref{sec:certificate}.  We first record its two consequences.

\begin{lemma}[Half-degree triangle moment]\label{lem:triangle-moment}
For every graphon of density \(p\in[2/3,1]\),
\begin{equation}
 I^2\geq p^2c_3T.
 \label{eq:triangle-moment}
\end{equation}
\end{lemma}

\begin{proof}
Integrating \eqref{eq:first-plane} and using \eqref{eq:means} gives
\[
 2I-rT\geq prc_3.
\]
Thus
\[
 I\geq\frac{\sqrt p}{2}(T+pc_3).
\]
Together with \(T\geq pc_3\), arithmetic--geometric mean yields
\[
 I\geq \sqrt p\sqrt{Tpc_3}=p\sqrt{c_3T}.
\]
Squaring proves \eqref{eq:triangle-moment}.  Equivalently, the exact surplus
before the last step is
\[
 \frac p4(T+pc_3)^2-p^2c_3T
 =\frac p4(T-pc_3)^2.
\]
\end{proof}

Let \(\kappa_4(x)\) denote the rooted \(K_4\) density at \(x\), and put
\[
 I_4=\int_\Omega d(x)^{1/2}\kappa_4(x)\dd\mu(x).
\]

\begin{lemma}[Half-degree adjacent-clique ratio]\label{lem:k4-ratio}
For every graphon of density \(p\in[2/3,1]\),
\begin{equation}
 I_4\geq c_4 I.
 \label{eq:k4-ratio}
\end{equation}
\end{lemma}

\begin{proof}
If \(d(x)>0\), use the neighbour probability measure
\[
 \dd\nu_x(y)=\frac{W(x,y)}{d(x)}\dd\mu(y).
\]
The edge density of the link graphon under \(\nu_x\) is
\[
 z_x=\frac{\tau(x)}{d(x)^2},
\]
whereas \(d(x)^{-3}\kappa_4(x)\) is its triangle density.  Goodman inside
the link therefore gives
\begin{equation}
 \kappa_4(x)\geq
 d(x)^3z_x(2z_x-1)
 =\frac{2\tau(x)^2}{d(x)}-d(x)\tau(x).
 \label{eq:link-goodman}
\end{equation}
The assertion is trivial on \(\{d=0\}\).  Multiply
\eqref{eq:link-goodman} by \(\sqrt d\), integrate, and then integrate
\eqref{eq:second-plane}.  The mean-zero correction terms disappear by
\eqref{eq:means}, leaving
\[
 I_4
 \geq
 2\int\frac{\tau^2}{\sqrt d}
 -\int d^{3/2}\tau
 \geq c_4\int\sqrt d\,\tau
 =c_4I.
\]
\end{proof}

\section{The page compression and the theorem}

For a spine triple \(\boldsymbol x=(x_1,x_2,x_3)\), write
\[
 A(\boldsymbol x)=
 W(x_1,x_2)W(x_1,x_3)W(x_2,x_3)
\]
and, for \(s\geq0\),
\[
 K_s(\boldsymbol x)=
 \int_\Omega
 d(y)^s\prod_{i=1}^3W(y,x_i)\dd\mu(y).
\]
Integrating the pendant leaf first gives the exact identity
\begin{equation}
 t(H_{178},W)=\int_{\Omega^3}A(\boldsymbol x)
 K_1(\boldsymbol x)K_0(\boldsymbol x)\dd\mu^3.
 \label{eq:page-identity}
\end{equation}
Cauchy--Schwarz in the page variable gives
\[
 K_{1/2}(\boldsymbol x)^2
 \leq K_0(\boldsymbol x)K_1(\boldsymbol x).
\]
Another weighted Cauchy--Schwarz inequality, now on the spine triple, gives
\begin{equation}
 t(H_{178},W)
 \geq\int A K_{1/2}^2
 \geq\frac{(\int A K_{1/2})^2}{\int A}
 =\frac{I_4^2}{T}.
 \label{eq:page-cauchy}
\end{equation}
There is no zero-denominator issue: \eqref{eq:goodman} gives
\(T\geq2/9\) on the required interval.

\begin{theorem}[Atlas 178]\label{thm:atlas178}
For every graphon \(W\) with edge density \(p\in[2/3,1]\),
\[
 t(H_{178},W)\geq p^2(2p-1)(3p-2)^2.
\]
\end{theorem}

\begin{proof}
Combine \eqref{eq:page-cauchy}, Lemma~\ref{lem:k4-ratio}, and
Lemma~\ref{lem:triangle-moment}:
\[
 t(H_{178},W)
 \geq\frac{I_4^2}{T}
 \geq c_4^2\frac{I^2}{T}
 \geq p^2c_3c_4^2.
\]
By \eqref{eq:target}, the last expression is the exact chromatic target.
\end{proof}

\section{Exact certificate for the rooted planes}
\label{sec:certificate}

We prove Lemma~\ref{lem:planes} entirely by rational polynomial
certificates.  Put
\[
 r=\frac{4+X}{5},\qquad p=r^2,\qquad \sqrt d=D,
 \qquad (X,D)\in[0,1]^2.
\]
Thus the certificate covers \(p\in[16/25,1]\), slightly more than the
required interval.  For fixed \(p,d\), the feasible polygon in the
\((a,t)\)-plane is cut out by the eight nonnegative slacks
\[
 a,\ d-a,\ a-d-p+1,\ t,\ t-2a+p,\
 t-a+d(1-d),\ a-t,\ d^2-t.
 \tag{C.1}
\]
Its possible vertices are
\begin{center}
\begin{tabular}{ccl}
\toprule
name & \(a\) & \(t\)\\
\midrule
\(A\)   & \(d\) & \(d^2\)\\
\(B0\)  & \(p/2\) & \(0\)\\
\(BL1\) & \(d^2-d+p\) & \(p-2d+2d^2\)\\
\(BL2\) & \(d+p-1\) & \(2d+p-2\)\\
\(BU1\) & \((d^2+p)/2\) & \(d^2\)\\
\(BU2\) & \(d\) & \(2d-p\)\\
\(BUp\) & \(p\) & \(p\)\\
\(C0\)  & \(d(1-d)\) & \(0\)\\
\(CL\)  & \(d+p-1\) & \(d^2+p-1\)\\
\bottomrule
\end{tabular}
\end{center}
Infeasible entries are discarded by (C.1).

For \eqref{eq:first-plane}, subtract the right side from the left.  This is
affine in \(a,t\), so it suffices to check the listed feasible vertices.
At \(BL1\), the residual contains the square
\((5D-X-4)^2\); its quotient is certified.  At \(BL2\), one constraint is
\(-(D^2-1)^2\geq0\), so \(D=1\), where the residual factors as
\[
 \frac{(X-1)^2(X+4)(2X+13)(2X+23)}{3125}.
\]
The exact conditional Bernstein trees are:
\begin{center}
\begin{tabular}{crrr}
\toprule
face & visited & accepted & infeasible\\
\midrule
\(A\)&115&29&29\\
\(B0\)&7&0&4\\
\(BL1\)&5&2&1\\
\(BU1\)&41&16&5\\
\(BU2\)&39&10&10\\
\(BUp\)&15&4&4\\
\(C0\)&13&6&1\\
\(CL\)&265&132&1\\
\midrule
total&500&199&55\\
\bottomrule
\end{tabular}
\end{center}
The only two equality-corner patches are also rational.  At \(A\), write
\(u=1-D,v=1-X\).  Feasibility gives
\[
 \frac{9}{25}v\leq1-p\leq(1-d)^2\leq4u^2,
 \qquad v\leq\frac{100}{9}u^2.
\]
The leading residual \(17u^2-v/5\) is therefore at least
\(133u^2/9\).  On \(u,v\leq1/64\), the absolute remainder is at most
\[
 \frac{6332030460179381}{4057816381784064}\,u^2
 <\frac{133}{9}u^2.
\]
At \(CL\), the leading quadratic is
\[
 17u^2-6uv+\frac35v^2.
\]
After subtracting \((u^2+v^2)/50\), its determinant is \(2121/2500\);
on \(u,v\leq1/2048\), the remainder is at most
\[
 \frac{697663143033}{53687091200000}(u^2+v^2)
 <\frac1{50}(u^2+v^2).
\]

For \eqref{eq:second-plane}, multiply its residual by \(D\).  The resulting
polynomial is
\begin{equation}
 Q=2t^2-(d^2+c_4d)t
 -3prc_3D(d-p)-2rc_3D(a-d^2).
 \tag{C.2}
\end{equation}
At \(D=0\), feasibility gives \(a=t=0\), and the unmultiplied inequality is
immediate.  For \(D>0\), \(Q\geq0\) is equivalent to
\eqref{eq:second-plane}.  The coefficient of \(a\) in (C.2) is
nonpositive.  Thus, for fixed \(d,t\), the minimum occurs on one of
\[
 a=d,\qquad a=\frac{t+p}{2},\qquad a=t+d(1-d).
\]
The first face collapses to \(A\).  On the other two faces, \(Q\) is a
convex quadratic in \(t\), so its minimum occurs at a listed endpoint or
at its stationary point.  This proves that the nine vertex checks and two
stationary checks below cover the whole feasible region.

\begin{center}
\begin{tabular}{crrr}
\toprule
face & visited & accepted & infeasible\\
\midrule
\(A\)&67&22&12\\
\(B0\)&5&0&3\\
\(BL1\)&5&2&1\\
\(BU1\)&11&3&3\\
\(BU2\)&25&6&7\\
\(BUp\)&1&1&0\\
\(C0\)&7&4&0\\
\(CL\)&277&138&1\\
\(B\)-stationary&7&0&4\\
\(C\)-stationary&1&1&0\\
\midrule
total&406&177&31\\
\bottomrule
\end{tabular}
\end{center}
Again, \(BL1\) and the \(C\)-stationary residual contain
\((5D-X-4)^2\), and \(BL2\) collapses to \(D=1\).  Its remaining boundary
factor is
\[
 \frac{(X-1)^2(X+4)(X+9)
 (6X^3+102X^2+503X+514)}{78125}.
\]
At the \(A\) corner, the leading quadratic is
\[
 18u^2-\frac{52}{5}uv+\frac{42}{25}v^2.
\]
After subtracting \((u^2+v^2)/10\), its determinant is \(621/500\); on
\(u,v\leq1/1024\), the remainder coefficient is
\[
 \frac{62190435867197668803}{900719925474099200000}<\frac1{10}.
\]

The \(CL\) face has a thin equality tube, which is treated directly instead
of forcing excessive bisection.  Put \(z=r-D\) and \(q=1-r^2\).  Its exact
residual is
\begin{align*}
 Q_{CL}
 &=\bigl(5r^6+22r^4-7r^2-2\bigr)z^2
   +2r(r^2+3)q^2z+2q^4+z^3R(r,z),\\
 R(r,z)
 &=-r(30r^4+31r^2-10)
 +(45r^4+23r^2-3)z\\
 &\quad-2r(21r^2+5)z^2+(25r^2+2)z^3
 -8rz^4+z^5.
\end{align*}
For \(9/10\leq r\leq1\), the first three terms are at least
\(\frac35(z^2+q^4)\); the two principal minors are certified by positive
one-variable Bernstein coefficients.  If \(|z|\leq1/100\), then
\[
 |z^3R(r,z)|
 \leq
 \frac{516552270801}{10^{12}}z^2
 <\frac35z^2.
\]
Thirty Bernstein boxes lie wholly inside this proved tube.  All other boxes
are settled by ordinary nonnegative Bernstein coefficients or by an
infeasible constraint.

Every subdivision, integer Bernstein tensor, corner coefficient, graph
identity, and target factorization is reproduced by
\[
 \texttt{codes/verify\_atlas178\_half\_degree\_weighted\_k4.py}.
\]

\section{Scope}

All integrations above are over an arbitrary probability space and an
arbitrary measurable symmetric kernel \(0\leq W\leq1\).  The proof uses
Tonelli--Fubini, Cauchy--Schwarz, Goodman in a normalized link, the rooted
feasible inequalities proved in Lemma~\ref{lem:feasible}, and the exact
finite polynomial certificates in Section~\ref{sec:certificate}.  It does
not assume that \(W\) is finite-step, regular, or bounded away from zero.
No limiting division at \(d=0\) is used.

\section{What the Lean formalization actually proves}
%%%%%%%%%%%%%%%%%%%%%%%%%%%%%%%%%%%%%%%%%%%%%%%%%%%%%%%%%%%%%%%%%%%%%%%%%%%%%%%

Theorem~\ref{thm:atlas178} is formalized in full, on \(p\in[2/3,1]\).  The
development is \texttt{lean/Taeyoung/Methods/Atlas178/}, the catalogue theorem
is \texttt{satisfiesLowerBound\_178}, and every declaration depends on exactly
\texttt{[propext, Classical.choice, Quot.sound]}.  Atlas \(178\) is
\texttt{verified}.

Sections~1, 2 and~4 are followed as written.
Section~\ref{sec:certificate} is not: none of it is used.

\subsection{Both planes are exact duals, not subdivision trees}

The two conditional Bernstein trees of \(500\) and \(406\) boxes, the
nine-vertex enumeration, the two equality-corner patches, and the \(Q_{CL}\)
tube are all replaced by three polynomial identities.  Writing the three
feasibility slacks of \eqref{eq:feasible} as
\[
 g_5=t-2a+p,\qquad g_6=t-a+d-d^{2},\qquad g_8=d^{2}-t,
\]
and putting \(\mathrm{res}_1\) for twice the residual of
\eqref{eq:first-plane} and \(\mathrm{res}_2\) for \(D\) times the residual of
\eqref{eq:second-plane}, the identities are
\begin{align*}
 \mathrm{res}_1&=\tfrac32 r\,g_5+\bigl(\tfrac72r-4D\bigr)g_8+\Phi_U
   &&\text{on }8D\leq7r,\\
 \mathrm{res}_1&=(5r-4D)g_5+(8D-7r)g_6+2(D-r)^{2}L
   &&\text{on }8D\geq7r,\\
 \mathrm{res}_2&=2Drc_3\,g_6+2(t-\theta)^{2}
   +\tfrac{D}{8}(D-r)^{2}M,
\end{align*}
with
\[
 \theta=\tfrac14 D(D+r)(D^{2}-Dr+4r^{2}-2),\qquad
 L=4D^{3}+6D^{2}r+8Dr^{2}-4D+4r^{3}-3r,
\]
\[
 \Phi_U=\tfrac12\bigl(8D^{5}-D^{4}r-24D^{2}r^{3}+6D^{2}r+16r^{5}-5r^{3}\bigr),
\]
\[
 M=-D^{5}-2D^{4}r-9D^{3}r^{2}+4D^{3}+40D^{2}r^{3}-20D^{2}r
   +80Dr^{4}-32Dr^{2}-4D+48r^{5}-24r^{3}.
\]
Every multiplier is nonnegative on its region: \(5r\geq4\geq4D\) because
\(r\geq4/5\), and \(7r/2\geq4D\) is the case hypothesis.  So the first plane
needs two cases and the second none.  The tangency at the balanced tripartite
graphon appears in both as an explicit \((D-r)^{2}\).

\subsection{The three leftover factors}

\(L\geq0\) on \(8D\geq7r\) is one \texttt{nlinarith}.  \(M\geq0\) on the whole
box is the four-term grouping
\[
 M=24r^{3}c_3+D(80r^{4}-41r^{2}-7)+20D^{2}rc_3
   +D\bigl(4D^{2}+9r^{2}+3-D^{4}-2D^{3}r-9D^{2}r^{2}\bigr).
\]
\(\Phi_U\geq0\) mixes degrees \(5\) and \(3\), so its degree-\(3\) part is first
traded against its degree-\(5\) part using \(25r^{2}\geq16\); what is left is
homogeneous, and in the band coordinate \(7r-8D\) it has positive coefficients
throughout:
\[
 32\,\Phi_U=H+r(25r^{2}-16)(5r^{2}-6D^{2}),\qquad
 H=128D^{5}-16D^{4}r-234D^{2}r^{3}+131r^{5},
\]
\[
 256\,H=2078r^{5}+2471(7r-8D)r^{4}
 +(7r-8D)^{2}\bigl(289r^{3}+1064r^{2}D+832rD^{2}+512D^{3}\bigr).
\]
No box subdivision, and no Bernstein coefficient, occurs anywhere in the
formalization.  The certificate interval is the same: \(r\geq4/5\) is
\(p\geq16/25\), the note's \(X\geq0\).

\subsection{The quotient at \(d=0\)}

The second plane is used in the divided form of \eqref{eq:second-plane}, since
that is what has mean-zero corrections.  Lean's \(x/0=0\) is exactly the note's
convention: at \(d=0\) feasibility gives \(\tau\leq d^{2}=0\), so both the
divided Goodman quotient \(\tau(2\tau-d^{2})/\sqrt d\) and the divided residual
vanish, while the right-hand side is \(-3p^{2}rc_3\leq0\).  The degenerate
fibre therefore needs no separate statement, only a separate two-line case.

\subsection{The page identity}

\eqref{eq:page-identity} is obtained by peeling all six coordinates and
collapsing the three page coordinates in turn; the leaf integral produces
\(d(x_3)\) and the two page integrals produce \(K_1\) and \(K_0\).  Both
Cauchy--Schwarz steps are the same lemma, once on \(\mu\) and once on the
assignment measure of \(\mathrm{Fin}\,3\), and the identification
\(\int A\,K_{1/2}=I_4\) is a single Fubini swap between the spine triple and
the page variable rather than three nested ones.

\end{document}
